% Insert after concurrent_base_slice_ja_v32.tex.
\subsection{共点三直線の五次消去障害と長さ六の厚み}
\label{subsec:concurrent-mixed-v32}

共点三直線では、基底切片の平方零方向に加えて、混合方程式から五次の消去障害が現れる。
以下の「五次」は非零 Poisson 点の近傍で線形方程式を消去した後の次数である。
原点における三十八個の Jacobi 式が全て二次であることと矛盾しない。

\begin{theorem}[共点三直線の全形式局所模型]
\label{thm:concurrent-mixed-v32}
$(A,B)$ の空でない Zariski 開集合で、次の主張が成り立つ。
この開集合は $(A,B)=(2,5)$ を含む。
\[
 c_0=xy(x+y),\qquad D_0=(Ax+By)\partial_z,\qquad [F]=0
\]
における全 Poisson スキームの完備局所環は
\[
 \boxed{
 \frac{\mathbb C[[q_1,\ldots,q_{18},a,b,u]]}
 {(ab,b^2,a^5u)}
 }
\]
と同型である。$u$ は混合核を三次式による $w$ の移動で割った一次元の残余方向である。
従って、この局所スキームには二つの被約十九次元枝があり、
その交叉は滑らかな十八次元である。
$u\ne0$ の枝は一般点でも非被約であり、その一般点の横断局所環は
\[
 \frac{K[[a,b]]}{(a^5,ab,b^2)}
\]
である。ここで $K$ はこの被約成分の関数体である。
その長さは六、基底は
\[
 1,a,a^2,a^3,a^4,b
\]
である。
\end{theorem}

\begin{proof}
補題~\ref{lem:concurrent-base-v32} の基底切片を用いる。
$\gamma=-(B-2A)/(A-2B)$ と置けば
\begin{align*}
 c&=xy(x+y)+a xyz+b(xz^2+\gamma yz^2),\\
 D_x&=\frac{a(B-2A)}3x-2B\gamma b z,\\
 D_y&=\frac{a(A-2B)}3y-2Abz,\\
 D_z&=Ax+By+\frac{a(A+B)}3z,
\end{align*}
かつ基底環の関係は $ab=b^2=0$ である。
混合写像を
\[
 M(a,b)=M_0+aM_1+bM_2
\]
と書く。源は Euler 商の三十五次元、標的は六次形式の二十八次元である。

\smallskip\noindent
\textbf{次数による消去の制御.}
単項式 $F_i=x^{n_1}y^{n_2}z^{n_3}$ に
$z$ 重量 $n_3-\mathbf 1_{i=3}$ を与え、標的の単項式には通常の $z$ 次数を与える。
このとき $M_0,M_1,M_2$ はそれぞれ重量を $0,1,2$ だけ上げる。
$(A,B)=(2,5)$ における $M_0$ のブロックは次の通りである。
\[
 \begin{array}{c|rrrrrrrr}
 z\text{ 重量}&-1&0&1&2&3&4&5&6\\\hline
 \text{源の次元}&5&10&8&6&4&2&0&0\\
 \text{標的の次元}&0&7&6&5&4&3&2&1\\
 \text{階数}&0&7&6&5&4&2&0&0
 \end{array}
\]
非零小行列式を固定すれば、この階数表は $(A,B)$ の非空開集合で保たれる。
源・標的の次元の上界から、ここで示した階数がさらに増えることはない。
従って $M_0$ の核は十一、余核は四次元であり、余核の重量は $4,5,5,6$ である。

三次式の Hamiltonian 移動の像は、この開集合で十次元である。
実際、$(2,5)$ では $p_0$ の非零係数の共通因子はなく、
$p_0\times\nabla g=hx$ なら $p_0$ との内積から $h=0$ となる。
次に $\nabla g=\lambda p_0$、$\lambda\in\mathbb C$ と書け、
curl を取ると $\lambda D_0=0$ なので $g=0$ である。
この単射性は開条件である。
三次式 $g$ に重量 $\deg_z(g)-1$ を与えると、Hamiltonian 移動も重量を保ち、
その次元は重量 $-1,0,1,2$ でそれぞれ $4,3,2,1$ である。
従って核をその像で割った一次元商は重量 $-1$ に集中する。
$(2,5)$ では $F_{\rm ex}=(0,0,x^4)$ をその生成元に取れる。

重量ごとに選んだ二十四次の可逆小行列を使って、二十四変数を消去する。
$b=0$ とした Schur 補完を $R(a)$ と書く。
その $a^j$ の係数は重量を $j$ 上げる。
また、Hamiltonian 移動は $a$ の全次数で実際の核に持ち上がる。
従って $R(a)$ の最初の非零係数は、$M_0$ の核を Hamiltonian 像で割った商を通る。
この商の重量は $-1$、余核の最小重量は $4$ だから、
\[
 R(a)\equiv0\pmod{a^5}.
\]
これは単一パラメータの標本計算ではなく、上の開集合全体に通用する次数の議論である。

$b$ の Schur 係数も Hamiltonian 像を殺す。
残る商の重量 $-1$ を二つ上げても余核には届かないので、その係数は零である。
基底環では $ab=b^2=0$ であるから、$b$ を含む他の項は存在しない。

\smallskip\noindent
\textbf{五次係数の非零性.}
残る必要条件は、$a^5$ の係数が非零であることだけである。
$(A,B)=(2,5)$ での有理数計算を明示する。
Euler 商の代表として $F_x$ の $x$ で割り切れる単項式を零とする。
重量順に固定した二十四次小行列の定数行列を $A_0$、$a$ 係数を $A_1$ とすると
\[
 \det A_0=-906992640000\ne0,\qquad (A_0^{-1}A_1)^5=0.
\]
従って小行列の逆行列は有限和
\[
 (A_0+aA_1)^{-1}
 =\sum_{j=0}^4(-aA_0^{-1}A_1)^jA_0^{-1}
\]
で正確に与えられる。
消去後の自由係数のうち、$F_z$ の二変数単項式の係数を $h_{ij}=[x^iy^j]F_z$ と書く。
Schur 補完の非零式は正確に
\[
 a^5 L\quad\text{と}\quad a^6 L,
 \qquad
 L=\frac{20h_{13}-100h_{22}+300h_{31}-700h_{40}}9
\]
だけである。これらはそれぞれ標的の $x^2z^4,xz^5$ の係数に現れる。
$F_{\rm ex}$ に対して $L=-700/9$ なので、五次係数は非零である。
この非零条件も $(A,B)$ の開条件であり、開集合をさらに縮めてよい。
付属の \texttt{verify\_concurrent\_mixed\_v32.py} は、
小行列の指定、有限逆行列、全 Schur 係数と $b$ 係数の消滅を標準ライブラリの
有理数演算だけで検証する。

\smallskip\noindent
\textbf{全局所環の認定.}
Hamiltonian 移動の十方向を実際の核の自由直和因子として選び、残る源の座標を $u$ とする。
混合写像は $F$ に線形なので、消去後の全関係は $u$ にある係数ベクトルを掛けたものである。
前段によりその係数は $a^5$ で割れ、少なくとも一成分は $a^5$ と単元の積である。
従って関係イデアルは正確に $(a^5u)$ となる。
基底の八自由変数と Hamiltonian の十自由変数を合わせれば表示した環を得る。

被約化の二つの素イデアルは $(a,b)$ と $(b,u)$ であり、
それぞれ十九次元、交叉は十八次元である。
$u$ が単元である一般点では関係が $(a^5,ab,b^2)$ となる。
上の六単項式が基底であるから長さ六を得る。
\end{proof}

\begin{remark}[一般点の厚みと低次検査]
この成分の一般点では、横断方向 $a$ の四次までの項は消えないが、五次で初めて零になる。
二次の初期イデアル $(ab,b^2)$ だけを見ても、この長さ六は分からない。
また、三角形点で現れた非被約構造は既約枝の境界の現象であったが、
ここでは一つの十九次元成分が一般点において既に非被約である。
本定理は $(A,B)$ の明記した非空開集合に関する主張であり、
例外的な係数値で同じ厚みが続くとは主張しない。
\end{remark}
